\documentclass[11pt,a4paper]{article}
\usepackage[T1]{fontenc}
\usepackage[utf8]{inputenc}
\usepackage{lmodern}
\usepackage[margin=27mm]{geometry}
\usepackage{amsmath,amssymb,amsthm,mathtools}
\usepackage{microtype,booktabs,array,longtable,enumitem}
\usepackage{xurl}
\usepackage[hidelinks]{hyperref}
\hypersetup{pdfauthor={Literature-report preparation with GPT-6 Astra Pro}}
\setlength{\emergencystretch}{3em}
\setlength{\parskip}{4pt}
\setlist{itemsep=3pt,topsep=5pt}
\allowdisplaybreaks[2]
\numberwithin{equation}{section}
\newtheorem{theorem}{Theorem}[section]
\newtheorem{proposition}[theorem]{Proposition}
\newtheorem{lemma}[theorem]{Lemma}
\newtheorem{corollary}[theorem]{Corollary}
\theoremstyle{definition}
\newtheorem{definition}[theorem]{Definition}
\theoremstyle{remark}
\newtheorem{remark}[theorem]{Remark}
\newcommand{\R}{\mathbb R}
\newcommand{\C}{\mathbb C}
\newcommand{\Q}{\mathbb Q}
\newcommand{\N}{\mathbb N}
\newcommand{\F}{\mathbb F}
\newcommand{\eps}{\varepsilon}
\newcommand{\dd}{\,\mathrm d}
\newcommand{\poly}{\operatorname{poly}}
\newcommand{\spec}{\operatorname{spec}}
\newcommand{\supp}{\operatorname{supp}}
\newcommand{\rank}{\operatorname{rank}}
\newcommand{\pin}{\texttt{7ef5c0cc5ca4b0e072b7a050fca09ab24bc6fa3a}}
\newcommand{\statusbox}[1]{\begin{quote}\small\noindent\textbf{Status and attribution.} #1\end{quote}}

\title{The announced Yang--Mills existence result\\
\large A literature and scope report for Phys-1, not a mass-gap solution}
\author{Prepared for \texttt{we-solve} from the public announcement and problem statement}
\date{26 October 2026}
\begin{document}
\maketitle
\section*{Abstract}
The October 26, 2026 announcement includes an existence claim for pure
four-dimensional Yang--Mills theory, within broader claims about interacting
gauge theories. This report identifies the matching Phys-1 catalogue problem,
separates existence from the additional mass-gap requirement, and records what
is publicly assessable. The physics proof artifacts are withheld according to
the announcement, so this document does not reconstruct their argument or claim
a solution. Elementary spectral examples explain why a mass gap cannot be
inferred from existence or from nonuniform gaps in approximations. The result
is a status report for editorial consideration, not new mathematical progress
on the Yang--Mills construction.

\statusbox{Claim category: Literature report. The physics manuscript and Lean
proof were not available in the public materials inspected. No independent
verification, full solution, or proved partial Yang--Mills theorem is claimed.
A detailed account of the unreleased proof cannot be supplied from these
sources without inventing mathematics. This report deliberately does not do so.}

\section{The matched catalogue problem}
The relevant entry is \emph{Yang--Mills Existence and Mass Gap}, Phys-1,
stable handle
\texttt{prb\_143040e1e85343ff92d7f92801ed2f08}
\cite{catalogue}. Its official specification is the Jaffe--Witten Clay problem
description, especially Sections 3--4 and printed page 6 \cite{clay}.

The target is a nontrivial four-dimensional continuum quantum Yang--Mills
theory for \emph{every compact simple gauge group} $G$. It includes suitable
local observable fields, positivity, covariance, locality and the required
short-distance structure. In its vacuum representation, the Hamiltonian $H$
is positive and self-adjoint, $H\Omega=0$, and the additional spectral target is
\begin{equation}\label{eq:gap}
 0<m:=\inf\bigl(\spec(H)\setminus\{0\}\bigr)<\infty.
\end{equation}
The spectrum and its gap must belong to the actual quantum theory, not a
classical field equation, a formal integral, or an unrelated operator. These
requirements must be met together; the displayed inequality is not a replacement
for the field-theoretic axioms.

For clarity, write the full target schematically as
\[
 \forall G:\quad \mathsf{Existence}(G)
       \ \text{and}\ \mathsf{Gap}(G),
\]
where $\mathsf{Existence}(G)$ includes the specified nontriviality and axiomatic
conditions. An existence announcement addresses only the first conjunct unless
a gap is separately stated and proved.

\section{What the public announcement establishes as a report}
OpenAI announces existence proofs for certain interacting gauge theories in four
dimensions, including QCD and a unified completion of the Standard Model. It
identifies pure Yang--Mills existence as a subproblem, says the approach builds
on Balaban's work, and reports completion of a Lean formalization of
Yang--Mills existence. The physics collaboration has elected to withhold the
AI-produced artifacts while preparing an exposition \cite{announcement}.

These are claims about what was obtained, not a public mathematical statement
with all hypotheses available for examination. The announcement does not provide
a mass-gap theorem. Its formalization statement concerns Yang--Mills existence;
it should not be enlarged into a claim that a mass gap, QCD, and the unified
theory have all been publicly formalized and verified.

\paragraph{Publicly accessible scope.}
For this report there is an announcement, an exact catalogue target and the
official Clay specification. There is no inspected public physics manuscript,
precise theorem signature, proof dependency list or reproducible formal-checking
configuration to cite as the proof of the announced construction. This is an
availability limit of the present source set, not evidence that the claimed
construction is false.

\paragraph{Quantifiers cannot be supplied by inference.}
Before matching even the existence part of Phys-1, a released theorem must show
which groups and parameters it covers, which continuum axioms it establishes,
and in what sense it is nontrivial. A general phrase such as ``pure Yang--Mills
exists'' cannot be silently expanded into all those exact quantifiers. This
report therefore records an announced existence advance relevant to Phys-1,
not a verified resolution of its existence component for every required group.

\section{Why existence does not imply a mass gap}
The distinction is mathematically substantive. Positivity and a vacuum do not
by themselves keep positive energies away from zero.

\begin{proposition}[An elementary spectral counterexample]\label{prop:gapless}
There is a bounded positive self-adjoint operator with a unique zero-energy
vector up to scalar multiples, nonzero excited states, and no positive gap.
\end{proposition}
\begin{proof}
On $\mathcal H=\C\Omega\oplus\ell^2(\N)$, with orthonormal excited basis
$e_1,e_2,\ldots$, define
\[
 H\Omega=0,\qquad He_j=\frac1j e_j.
\]
It is a bounded diagonal self-adjoint operator. For
$\psi=a\Omega+\sum_j b_je_j$,
\[
 \langle\psi,H\psi\rangle=\sum_{j\ge1}\frac{|b_j|^2}{j}\ge0,
\]
and $H\psi=0$ implies every $b_j=0$. Its spectrum is the closure of its
diagonal entries, namely $\{0\}\cup\{1/j:j\ge1\}$. Hence
\[
 \inf\bigl(\spec(H)\setminus\{0\}\bigr)=0,
\]
even though the vacuum is unique and there are nonzero excitations.
\end{proof}

This is an elementary operator example, \emph{not} a Yang--Mills theory or a
counterexample to the Clay problem. Its only purpose is to invalidate an
inference based solely on existence of a positive-energy Hilbert-space
representation and a vacuum. Extra field-theoretic information would have to
supply the desired gap.

\subsection{Positive gaps at every cutoff can still disappear}
\begin{proposition}[Nonuniform gaps do not survive even norm convergence]
There are bounded positive self-adjoint operators $H_N$ with the same unique
vacuum and with strictly positive gap for every $N$, converging in operator
norm to the gapless operator in Proposition~\ref{prop:gapless}.
\end{proposition}
\begin{proof}
On the same space set
\[
 H_N\Omega=0,\qquad H_Ne_j=\max\{1/j,1/N\}\,e_j.
\]
Each $H_N$ has gap $1/N$. Moreover,
\[
 \|H_N-H\|=\sup_{j\ge1}\bigl(\max\{1/j,1/N\}-1/j\bigr)=1/N.
\]
Thus $H_N\to H$ in norm while the positive gaps tend to zero. The limit is
nontrivial and has a unique vacuum, but is gapless.
\end{proof}

Again this is an illustrative spectral calculation, not a lattice gauge model.
It shows that a construction using approximations needs a lower bound on the
physical gap that does not vanish in the required limits, together with a valid
transfer of that bound to the continuum theory. Merely establishing a positive
gap for each approximation is insufficient.

\section{The proof obligations a released construction must resolve}
The following is an audit map derived from the catalogue target, not a summary
of hidden steps in OpenAI's unavailable proof.

\paragraph{Exact mathematical object.}
A theorem should specify the group, spacetime, observables, parameters and
representation being constructed. A QCD theorem with matter is not automatically
a pure-gauge theorem, and a theorem for one selected group does not automatically
meet the universal group quantifier. Any transfer between these objects needs
its own argument.

\paragraph{Continuum and nontriviality.}
A cutoff measure or a regulated Hamiltonian may be a useful starting object.
Its relation to a continuum theory and its nontriviality must still be proved.
A limiting theory in which all relevant observables become trivial would not
meet the target. The needed compactness, renormalization and limit arguments
cannot be filled in from the name of a constructive approach.

\paragraph{Axioms and physical observables.}
A proof must establish the required positivity, covariance and locality for the
actual observable system, and match the short-distance requirements. For a
Euclidean route, ordinary positivity of a probability measure must not be
confused with the reflection-positivity property needed by the chosen
reconstruction framework. A direct Hamiltonian route has its own operator and
representation obligations. These are alternatives to be identified in a
released proof, not assumptions this report makes about its method.

\paragraph{The full excitation space.}
A bound for one selected correlation function or one special channel does not,
without a theorem linking it to the whole vacuum-orthogonal space, determine
$\inf(\spec(H)\setminus\{0\})$. The spectral conclusion must concern the
physical Hilbert space and all required nonvacuum states. The finite upper
bound in \eqref{eq:gap} also rules out satisfying the request merely by leaving
no finite-energy excitation at all.

\paragraph{Uniformity of any limiting gap estimate.}
If regularization and volume limits are used, the relevant physical lower bound
must survive those limits. The preceding example explains why positivity at
each intermediate stage is not enough. This is a logical requirement for an
argument of that form, not a claim that OpenAI uses or lacks such an estimate.

\paragraph{Formalization alignment.}
A public formal theorem would need explicit definitions, an auditable dependency
closure and a reproducible validation procedure. One must compare the formal
conclusion with the continuum and spectral requirements above. A formal proof
of an existence assertion would not also prove an omitted mass-gap assertion.
The announcement does not make those unavailable artifacts inspectable.

\section{Relation to the earlier Phys-1 contribution}
The existing Phys-1-A1 exploration distinguishes a finite-cutoff construction
from a continuum theory and discusses conditions under which correlation decay
would imply a gap \cite{catalogue}. This report does not re-label those standard
conditional deductions as new discoveries. Its contribution is the catalogue
scope assessment of the new announcement and the separation between reportable
existence progress and an unreported full solution.

The elementary examples above are included for clarity, not as novel partial
progress on Yang--Mills theory. Under the site's significance criterion,
editors may prefer to record this document as a status update rather than a
research attempt. It should not be promoted to ``Full solution'' or to a proved
``Partial result'' merely because the underlying announcement is important.

\section{Exact outcome and what remains unavailable}
The appropriate present outcome is:
\begin{quote}
An announced, externally attributed existence advance relevant to Phys-1 is
recorded. The public material inspected does not provide the construction for
review, does not establish its exact universal-group match here, and does not
supply the required mass-gap theorem. No resolution of the full catalogue
problem is claimed by this submission.
\end{quote}

When source artifacts become public, the report can be replaced or extended by
a theorem-by-theorem exposition with precise hypotheses, source versions and a
verification record. Until then, fabricating an auxiliary equation, convergence
estimate, mass-gap proof or Lean theorem name would misrepresent the available
evidence. This document intentionally contains none of those invented details.

\section{Submission disclosures}
\paragraph{Problem handle.} \texttt{prb\_143040e1e85343ff92d7f92801ed2f08}.
\paragraph{Claim.} Literature report.
\paragraph{Exact scope.} Reports the October 26 announcement and identifies its relation
to the existence component of Phys-1. It does not supply an unreleased proof,
establish the exact universal-group scope, or prove a mass gap. The spectral
examples are explanatory and are not claimed as new Yang--Mills results.
Editorial status update / literature report only.
\paragraph{Who did the work.} LLM only (people only prompted or submitted).
\paragraph{Models used in preparing this report.} OpenAI GPT 6 Astra Pro.
\paragraph{Thinking level / time.} Unknown (not available).
\paragraph{Per-model versions, thinking and roles.}
GPT-6 Astra Pro; exact version: unknown; thinking setting: unknown;
actual compute budget/duration: unknown; role: source retrieval, mathematical
exposition, catalogue-scope comparison, and preparation of this report.
No second model or independent human proof reviewer was used in this preparation.
\paragraph{People and roles.}
Submitted as a standard We-Solve addition. Human roles are prompting and submission; no human proof-discovery or expert proof-verification role is claimed. Authenticated submitter identity and declarations are retained in the contribution metadata.
\paragraph{Original research provenance.} The announced research is attributed to OpenAI. The independent academic collaboration is not represented here as having authored, endorsed, or validated the unavailable proof.
The internal model names, versions, thinking settings and budgets behind those
source results are not established by the inspected material and are recorded as
unknown. They are external source provenance, not additional models invoked to
prepare this report.
\paragraph{Eligibility and verification evidence.}
Astra Pro is an example model in the inspected submission form. This is not
independent verification of the source research. The exact checks and omissions
are stated in the body of the report.
\paragraph{Rights and declarations.}
This is newly written exposition with attribution and links, not a redistribution
of the source manuscripts or their code. New exposition is intended for the
site's CC BY 4.0 terms. The submitting human explicitly confirmed publication rights, public release, and complete people/model disclosure on 26 October 2026; cited third-party materials retain their own terms. No source papers, PDFs or executable code are redistributed.
\begin{thebibliography}{9}
\bibitem{catalogue} \emph{we-solve}, Phys-1, Yang--Mills Existence and Mass Gap, and its existing contribution, inspected 26 October 2026. Used: exact target, stable handle and distinction from previous conditional observations.
\url{https://we-solve.github.io/problems/yang-mills-existence-and-mass-gap-143040e1/}.
\bibitem{clay} Arthur Jaffe and Edward Witten, \emph{Quantum Yang--Mills Theory}, official Clay Mathematics Institute problem description, Sections 3--4, especially printed page 6. Used: universal gauge-group target, axiomatic conditions and mass-gap definition. The relevant PDF page was inspected.
\url{https://www.claymath.org/wp-content/uploads/2022/06/yangmills.pdf}.
\bibitem{announcement} OpenAI, \emph{Working with academic communities to build scientific understanding}, 26 October 2026. Used: the announced existence scope, reported Yang--Mills formalization and decision not to release the physics artifacts.
\url{https://openai.com/index/building-scientific-understanding-with-academic-communities/}.
\end{thebibliography}
\end{document}
